\justifying
\NSFCBodyText

\subsubsection{已具备的实验条件}
申请人已确认可稳定使用 NVIDIA A40 共8张、NVIDIA GeForce RTX 3090 Ti 共6张，并具备 Linux、Python、深度学习框架、Git 与 LaTeX 等软件工具链。这些条件适合中小规模指令模型推理、轻量压缩器训练及分批重复评测。实验根据显存与并发条件配置，不假定不同型号显卡可以组成同质高速互联集群，也不把总卡数直接换算为大模型训练能力。

任务状态保真研究主要需要轨迹记录、可复位环境和少量标注；有效恢复研究需要原文档案、对象索引和版本查询；调度研究需要统一日志和调用成本记录。上述能力可在本地软件环境逐步实现。AppWorld、WebArena 等公开资源作为拟获取并按许可证使用的实验来源，其可公开获取不等于项目已完成本地部署或获得全部数据授权。

\subsubsection{尚缺少的实验条件及拟解决的途径}
本项目尚缺共同条件下的智能体先导轨迹、动作依赖标注和完整对照日志，计划在启动初期建立。模型、任务与依赖包版本在评测前固定；项目级数据和代码的形成过程保留来源与变更记录，以支持后续独立复跑。团队的实际成员、分工和投入安排按已登记事实缺口补齐后落实。

企业侧业务验证尚需\textbf{【待补 F-KS-03：数据授权范围、业务接口和协作负责人】}。基金通知中的技术、算力和脱敏数据支持属于合作安排，不能等同于本申请已获得固定配额。核心算法和验证依托现有算力及公开环境完成；额外企业算力、学院公共资源或云端模型仅作为落实后的补充条件，不预写未获授权的设备与配额。

\subsubsection{利用研究基地的计划与落实情况}
项目依托申请人已确认的本地计算条件开展，不以国家实验室、全国重点实验室或部门重点实验室的未落实资源作为前提。需要使用学院公共算力时，按实际申请结果安排任务；若补充资源不能及时获得，优先减少原生大模型复现与扩展场景规模，保留核心交叉实验、强组合对照和多次重复，避免牺牲主要证据链。
